\section{Related Work}
\label{sec:related}

We survey three areas relevant to SA-based code quality prediction: static analysis for LLM code, self-repair and iterative refinement, and code generation evaluation. Our positioning: prior work uses SA for feedback (correction) rather than prediction (correlation)---we fill this gap.

\subsection{Static Analysis for LLM-Generated Code}

Recent studies integrate static analysis into LLM code generation pipelines, but as feedback rather than predictors. \citet{Blyth2025Static} demonstrate an iterative feedback loop where Bandit and Pylint errors are fed back to the LLM for self-correction. Their approach reduces security vulnerabilities from $>$40\% to 13\% and readability violations from $>$80\% to 11\% within 10 iterations on HumanEval/MBPP. However, they measure improvement magnitude, not predictive correlation---whether SA scores before correction predict which samples will pass tests.

AutoSafeCoder~\citep{Nunez2024AutoSafeCoder} combines multi-agent SA with fuzz testing, achieving 13\% reduction in code vulnerabilities. CodeQUEST~\citep{Liu2025CodeQUEST} uses GPT-4o to iteratively improve code quality as measured by Pylint Score, Radon Maintainability, and Bandit logs, reporting 52.6\% mean improvement and ``meaningful correlation'' between SA metrics and quality. Yet neither study quantifies this correlation---no $r$-values or regression coefficients appear.

STALL+~\citep{Liu2024STALL} integrates static analysis at the prompting phase for repository-level code completion, finding SA integration performs best when combined with RAG. Their focus is on improving generation quality through SA-informed prompts, not on using SA metrics as standalone predictors of correctness.

Our work differs fundamentally: we measure correlation directly ($r=0.87$ for pylint-correctness), enabling SA-based rejection sampling without any LLM feedback loop.

\subsection{Self-Repair and Iterative Refinement}

The self-repair paradigm feeds execution feedback---including SA errors---back to LLMs for correction. \citet{Olausson2024SelfRepair} systematically study self-repair effectiveness, finding it improves pass rates but is not universally effective across models. \citet{Arimbur2026Tries} shows self-repair improves HumanEval pass rates by +4.9 to +17.1 percentage points, with most gains occurring in the first 2 iterations.

CodeCoR~\citep{Pan2025CodeCoR} achieves 77.13\% average Pass@1 on HumanEval/MBPP using multi-agent pruning and test case generation. INTERVENOR~\citep{NEUIR2024INTERVENOR} employs Code Teacher and Code Learner agents with compiler feedback. RePair~\citep{TnTWoW2024RePair} uses process-based feedback with a reward model as critic.

These approaches share a common assumption: SA feedback is useful for correction. Our work tests a logically prior question: does SA signal predict correctness in the first place? If correlation is weak, iterative feedback may work through mechanisms other than SA signal. Our finding that $r=0.87$ validates the assumption underlying these approaches.

\subsection{Code Generation Evaluation}

Benchmark infrastructure underpins code generation research. \citet{Chen2021Evaluating} introduced HumanEval with 164 hand-crafted problems and the pass@$k$ metric. \citet{Austin2021Program} created MBPP with 974 problems (427 sanitized). \citet{Liu2023EvalPlus} extended these with EvalPlus, adding 80$\times$ more test cases to HumanEval+ and revealing that 19--29\% of previously ``correct'' code actually fails under rigorous testing.

\citet{Yetistiren2023Evaluating} evaluate Copilot, CodeWhisperer, and ChatGPT on HumanEval across correctness, security, reliability, and maintainability dimensions, using Radon, Bandit, and Pylint as quality metrics. However, they report quality scores and correctness separately---no correlation analysis between the two.

We build on this evaluation infrastructure while filling the correlation gap: using HumanEval+MBPP as ground truth, we measure how well SA metrics predict pass@1 outcomes.

\begin{table}[t]
\centering
\caption{Comparison of SA usage in prior work. We are the first to quantify SA-correctness correlation.}
\label{tab:related}
\begin{tabular}{@{}llll@{}}
\toprule
Work & SA Usage & Correlation Measured & $r$-Value \\
\midrule
\citet{Blyth2025Static} & Feedback loop & No & --- \\
CodeQUEST~\citep{Liu2025CodeQUEST} & Iterative improvement & Qualitative & --- \\
AutoSafeCoder~\citep{Nunez2024AutoSafeCoder} & Multi-agent feedback & No & --- \\
\citet{Yetistiren2023Evaluating} & Quality measurement & No & --- \\
\textbf{Ours} & \textbf{Prediction} & \textbf{Yes} & $\mathbf{r=0.87}$ \\
\bottomrule
\end{tabular}
\end{table}
